\chapter{Medical Emergencies}\label{ch:medical}

\section{Heart Attack (Myocardial Infarction)}

A heart attack occurs when blood flow to part of the heart muscle is blocked, often by a blood clot in a coronary artery. Prompt recognition and response are critical for survival and reducing heart damage.

\subsection{Signs and Symptoms}

\begin{itemize}[leftmargin=*]
  \item Chest pain or discomfort: Pressure, squeezing, fullness, or pain in center or left side of chest lasting more than a few minutes or coming and going
  \item Pain radiating to arm (especially left arm), back, neck, jaw, or stomach
  \item Shortness of breath with or without chest discomfort
  \item Breaking out in a cold sweat
  \item Nausea or vomiting
  \item Lightheadedness or sudden dizziness
  \item Unusual fatigue (particularly in women)
\end{itemize}

\noindent\textbf{\textcolor{red}{WARNING}}: \textit{Symptoms vary by gender}: Women are more likely to experience shortness of breath, nausea/vomiting, and back/jaw pain. Both men and women can experience unusual fatigue. Older adults and people with diabetes may have milder symptoms or no chest pain at all.

\subsection{Immediate Actions}

\begin{enumerate}[leftmargin=*]
\item Have the person stop all activity and rest comfortably -- sitting or semi-reclining position is usually best
\item Call emergency services immediately (dial 112 or 108 in India). Do not drive the person yourself unless absolutely necessary
\item If prescribed nitroglycerin is available and the person is conscious, assist them in taking it as directed (usually one tablet under the tongue every 5 minutes up to three doses)
\item Ask if they carry aspirin; if yes and they are not allergic (and have not been advised by their doctor against it), chew one adult aspirin (325 mg) or four low-dose aspirins (81 mg each) -- chewing absorbs faster \cite{aha_arc_2024}
\end{enumerate}
\noindent\textbf{\textcolor{red}{NOTE}}: \textit{ASPIRIN WARNING}: Only give aspirin if the person has NO allergy to it, has NOT been advised by a doctor against taking it, and you are reasonably sure they are having a heart attack (not a bleeding condition like stroke). Do NOT give aspirin to children or teenagers.

\begin{enumerate}[leftmargin=*]
\item Monitor breathing continuously. If they become unresponsive and not breathing normally, begin CPR immediately
\end{enumerate}
\noindent\textbf{\textcolor{red}{DO NOT}}: Give anything by mouth (food, drink, other medications) unless specifically prescribed; let the person walk around to ``walk off'' the discomfort; assume it's just indigestion without medical evaluation; drive the person to the hospital yourself.

\section{Stroke}

A stroke occurs when blood supply to part of the brain is interrupted or reduced, preventing brain tissue from getting oxygen and nutrients. Brain cells begin to die within minutes.

\subsection{Recognition -- The FAST Mnemonic}

\begin{center}
\begin{tabular}{@{}ll@{}}
\textbf{F} -- Face & Ask the person to smile. Is one side of the face drooping? \\
\textbf{A -- Arm} & Ask the person to raise both arms. Does one arm drift downward? \\
\textbf{S -- Speech} & Ask the person to repeat a simple phrase. Is speech slurred or strange? \\
\textbf{T -- Time} & If any of these signs are present, call emergency services immediately! \\
\end{tabular}
\end{center}

Additional stroke signs include:

\begin{itemize}[leftmargin=*]
  \item Sudden confusion or trouble understanding speech
  \item Sudden trouble seeing in one or both eyes
  \item Sudden severe headache with no known cause ("worst headache of my life")
  \item Sudden trouble walking, dizziness, loss of balance or coordination
\end{itemize}

\subsection{Management}

\begin{enumerate}[leftmargin=*]
\item Call emergency services immediately -- time is brain
\item Note the exact time when symptoms first appeared. This is critical information for medical treatment decisions
\item Keep the person comfortable and calm. Loosen tight clothing
\item Do NOT give food, drink, or medication (risk of aspiration if consciousness changes)
\item If the person becomes unconscious but is breathing, place them in the recovery position
\end{enumerate}
\noindent\textbf{\textcolor{red}{WARNING}}: Even if symptoms resolve quickly, this could be a transient ischemic attack (TIA or "mini-stroke"), which requires immediate medical evaluation as it often precedes a full stroke. Always seek emergency care for suspected stroke.

\section{Seizures (Epileptic Seizure)}

A seizure is a disturbance in normal electrical activity in the brain that can cause changes in behavior, movements, feelings, or levels of consciousness.

\subsection{Types of Seizures}

\begin{description}[leftmargin=!]
  \item[Generalized Tonic-Clonic (Grand Mal)] Loss of consciousness, body stiffens (tonic phase), followed by rhythmic jerking (clonic phase)
  \item[Absence (Petit Mal)] Brief staring spells (seconds), mostly in children
  \item[Partial/Focal] May involve localized jerking, altered awareness, or sensory changes without full loss of consciousness
\end{description}

\subsection{During a Seizure}

\begin{enumerate}[leftmargin=*]
\item Stay calm and time the seizure
\item Protect the person from injury by moving objects away and placing something soft under their head
\item Turn the person onto their side if possible to help maintain airway and allow fluids to drain
\item Loosen tight clothing around the neck
\end{enumerate}
\noindent\textbf{\textcolor{red}{NEVER}}: Put anything in the person's mouth (can cause broken teeth or jaw injury); try to force the person awake during the seizure; offer water or pills until fully alert and able to swallow safely.

\subsection{After the Seizure (Post-Ictal Phase)}

\begin{itemize}[leftmargin=*]
  \item Stay with the person until they are fully awake and oriented
  \item Offer reassurance and explain what happened
  \item Allow them to rest; some feel tired or confused afterward
  \item Do not force them to eat or drink until fully conscious and able to swallow safely
  \item Check for injuries that occurred during the seizure
\end{itemize}

\noindent\textbf{\textcolor{red}{When to Call Emergency Services}}: First-time seizure; seizure lasts longer than 5 minutes; second seizure occurs before the person recovers; difficulty breathing after seizure; person is injured; person does not regain consciousness; person has diabetes or is pregnant; seizure occurs in water.

\noindent\textbf{\textcolor{green}{PEARL}}: \textit{Seizure First Aid Summary}
\begin{enumerate}
  \item Protect from injury (move things away, pad head)
  \item Turn person onto their side
  \item Time the seizure
  \item Do not put anything in the mouth
  \item Do not restrain
  \item Stay with person until fully recovered
\end{enumerate}

\section{Diabetic Emergencies}

Both high blood sugar (hyperglycemia) and low blood sugar (hypoglycemia) are potential emergencies in people with diabetes.

\subsection{Hypoglycemia (Low Blood Sugar)}

Occurs when blood glucose falls below normal levels. Common causes: missed meals, excessive insulin, increased physical activity.

\subsubsection{Symptoms}

\begin{itemize}[leftmargin=*]
  \item Shakiness, sweating, trembling
  \item Rapid heartbeat or palpitations
  \item Hunger
  \item Dizziness, lightheadedness
  \item Confusion, irritability, anxiety
  \item Blurred vision
  \item Weakness, fatigue
  \item In severe cases: seizures, loss of consciousness
\end{itemize}

\subsubsection{Management (Conscious Person)}

\begin{enumerate}[leftmargin=*]
\item Administer fast-acting carbohydrates: 15 grams of simple sugar (glucose tablets, juice, regular soda, honey)
\item Wait 15 minutes and check blood sugar if possible
\item Repeat if still low
\item Once blood sugar is stable, provide a snack or meal containing protein and complex carbohydrates
\end{enumerate}
\noindent\textbf{\textcolor{red}{UNCONSCIOUS PERSON}}: Do NOT attempt to give anything by mouth (choking hazard). Apply glucose gel/tube inside the cheek if available and trained to use. Call emergency services immediately. Administer glucagon injection if prescribed and trained. Do NOT give food or drink until fully conscious and able to swallow.

\subsection{Hyperglycemia (High Blood Sugar)}

Occurs when there is too much glucose in the blood. Can progress to diabetic ketoacidosis (DKA) in type 1 diabetes or hyperosmolar hyperglycemic state (HHS) in type 2 diabetes.

\subsubsection{Symptoms}

\begin{itemize}[leftmargin=*]
  \item Extreme thirst
  \item Frequent urination
  \item Blurred vision
  \item Fatigue, weakness
  \item Headache
  \item Dry mouth and skin
  \item Fruity-smelling breath (DKA warning sign)
\end{itemize}

\subsubsection{Management}

\begin{enumerate}[leftmargin=*]
  \item Check blood sugar if possible
  \item Help the person take prescribed insulin only if they are conscious, alert, and able to swallow safely
  \item Encourage drinking water (if conscious, alert, and able to swallow safely — do NOT force fluids)
  \item Seek medical attention immediately if blood sugar remains extremely high (>300 mg/dL in adults), or if there are signs of ketones, fruity breath, confusion, or difficulty breathing
  \item If the person is unconscious, confused, or unable to swallow safely, DO NOT give food, drink, or insulin — seek emergency medical help immediately
\end{enumerate}
\noindent\textbf{\textcolor{red}{RECOGNIZING DKA}}: Look for rapid deep breathing (Kussmaul breathing), abdominal pain, nausea/vomiting, and fruity breath odor. This is a life-threatening emergency requiring immediate hospital care.


\section{Bee and Insect Stings}

Bees, wasps, hornets, and ants can sting, causing pain, redness, and swelling at the sting site. Most stings produce only local symptoms, but some people may develop a severe allergic reaction (anaphylaxis).

\subsection{Local Reaction (Most Common)}

These symptoms affect only the area around the sting and typically improve within hours to days.

\begin{itemize}[leftmargin=*]
  \item Immediate sharp pain or burning at the site
  \item Red welt at the sting location
  \item Localized swelling and itching
\end{itemize}

\subsubsection{Management}

\begin{enumerate}[leftmargin=*]
  \item Wash the area with soap and water to reduce infection risk
  \item Remove the stinger if present by scraping it sideways with a fingernail or credit card — DO NOT squeeze or pinch (this can inject more venom)
  \item Apply a cold pack or ice wrapped in cloth for 10-15 minutes at a time to reduce pain and swelling
  \item Elevate the affected limb if possible
  \item Use over-the-counter pain relievers (acetaminophen or ibuprofen) for pain
  \item Apply calamine lotion or hydrocortisone cream to reduce itching
  \item Take an oral antihistamine (e.g., diphenhydramine/Benadryl) if available and appropriate for the person's age and health
  \item Monitor for signs of allergic reaction that spread beyond the sting site (see below)
\end{enumerate}
\noindent\textbf{Note}: A large local reaction (increasing redness and swelling over several days, up to 4 inches in diameter) is different from an allergic reaction and usually resolves on its own. If swelling continues to worsen after 48 hours or shows signs of infection (pus, increasing warmth, red streaks), seek medical care.

\subsection{Systemic Allergic Reaction (Anaphylaxis)}

Some people experience a whole-body allergic reaction that can become life-threatening within minutes. This requires immediate action even if the person has never had a severe reaction before.

\textbf{Symptoms may include}: Skin hives or swelling away from the sting site; difficulty breathing, wheezing, or throat tightness; nausea, vomiting, or dizziness; rapid heartbeat; feeling of impending doom.

\textbf{If any systemic symptoms develop, treat as anaphylaxis immediately}: See the Anaphylaxis section below for emergency management including epinephrine administration and calling emergency services.


\section{Anaphylaxis (Severe Allergic Reaction)}

Anaphylaxis is a life-threatening systemic allergic reaction that can occur rapidly after exposure to an allergen. Common triggers: foods (peanuts, tree nuts, shellfish, eggs, milk), insect stings, medications, latex.

\subsection{Symptoms}

\begin{itemize}[leftmargin=*]
  \item Skin: Hives, itching, flushing, swelling (especially lips, tongue, face, throat)
  \item Respiratory: Wheezing, difficulty breathing, stridor (high-pitched sound when inhaling), hoarseness, coughing
  \item Cardiovascular: Rapid or weak pulse, pale/blue skin, dizziness, fainting, low blood pressure
  \item Gastrointestinal: Nausea, vomiting, abdominal cramps, diarrhea
  \item Others: Sense of impending doom, anxiety
\end{itemize}

\noindent\textbf{\textcolor{red}{EMERGENCY}}: Anaphylaxis can progress rapidly within minutes. Delay in epinephrine administration increases risk of death.

\subsection{Management}

\begin{enumerate}[leftmargin=*]
\item Identify and remove the allergen if possible (e.g., sting stinger scraped out sideways with fingernail)
\item Assist the person in using their prescribed epinephrine auto-injector (EpiPen, Auvi-Q, etc.) if available
\item Call emergency services immediately
\item Lie the person flat with legs elevated (unless breathing difficulty makes this uncomfortable, then allow semi-upright position)
\item Monitor breathing and circulation; prepare to perform CPR if needed
\item A second dose of epinephrine may be given after 5-15 minutes if symptoms persist or worsen and another auto-injector is available
\end{enumerate}
\noindent\textbf{\textcolor{green}{NOTE}}: \textit{Epinephrine Administration}: Inject into the outer thigh muscle through clothing if necessary. Hold in place for several seconds to ensure full dose delivery. Massage the injection site gently afterwards.

\noindent\textbf{\textcolor{green}{PEARL}}: \textit{Allergy Action Plan Reminder}

Always ask individuals with known severe allergies if they carry an epinephrine auto-injector and how you should assist if needed. Many people wear medical identification indicating their allergies and emergency instructions.

\section{Asthma Attack}

An asthma attack involves narrowing of the airways due to inflammation and muscle constriction, causing breathing difficulties.

\subsection{Symptoms}

\begin{itemize}[leftmargin=*]
  \item Difficulty breathing or shortness of breath
  \item Wheezing (whistling sound when breathing)
  \item Persistent coughing
  \item Chest tightness or pain
  \item Anxiety, panic from difficulty breathing
  \item Pale, sweaty, or bluish skin (severe attack)
\end{itemize}

\subsection{Management}

\begin{enumerate}[leftmargin=*]
\item Help the person remain calm and sit upright (this helps air movement better than lying down)
\item Assist in using their rescue inhaler (usually albuterol/salbutamol): shake, prime if new, take one puff and wait one minute, repeat as prescribed (typically 1-2 puffs every 4-6 hours as needed)
\item Call emergency services if: no improvement after 2-3 puffs over 10 minutes; speaking in words instead of sentences; lips/fingernails turning blue; feeling increasingly worse; never been diagnosed before; symptoms recur after initial improvement
\end{enumerate}
\noindent\textbf{\textcolor{red}{WARNING}}: Never give anyone else's asthma inhaler unless specifically authorized. Do not slap their back during an attack (this can make it worse).

\section{Other Medical Emergencies}

\subsection{Unexplained Loss of Consciousness (Syncope/Passing Out)}

Syncope is a brief, temporary loss of consciousness caused by a temporary drop in blood flow to the brain. Most causes are benign, but some are dangerous.

\subsubsection{Common Causes}

\begin{itemize}[leftmargin=*]
  \item \textbf{Vasovagal (simple) fainting}: triggered by pain, fear, anxiety, the sight of blood, or standing for long periods
  \item \textbf{Orthostatic (postural) hypotension}: a sudden drop in blood pressure when standing up quickly, especially in older adults or after dehydration
  \item \textbf{Dehydration or low blood sugar}
  \item \textbf{Dangerous causes}: heart rhythm problems, heart attack, internal bleeding, or stroke -- these require immediate medical care
\end{itemize}

\subsubsection{Immediate Management}

\begin{enumerate}[leftmargin=*]
  \item Lay the person flat on their back
  \item Raising the legs above heart level may be considered ONLY if the cause is non-traumatic (e.g., simple fainting) and it does not cause pain or worsen symptoms
  \item Loosen tight clothing
  \item Check breathing and pulse; if not breathing normally, begin CPR
  \item If conscious and recovering, allow them to rest and sit up gradually -- do not let them stand quickly
  \item Offer small sips of water once fully alert (if the cause may be dehydration) or a source of quick sugar if low blood sugar is suspected
  \item Seek medical evaluation for any unexplained syncope
\end{enumerate}
\subsubsection{Red Flags -- Call Emergency Services}

\begin{itemize}[leftmargin=*]
  \item Fainting during or right after exertion, exercise, or physical activity
  \item Fainting without any warning signs (no nausea, sweating, or lightheadedness beforehand)
  \item Fainting with palpitations (racing or fluttering heartbeat) or chest pain
  \item The person does not recover quickly or is confused after waking
  \item Repeated episodes, fainting during pregnancy, or any known heart condition
  \item The person is injured from the fall, or is taking blood-thinning medication
\end{itemize}

\quicknote{The person who wakes up quickly and feels well after a simple fainting spell (stress, heat, or standing too long) is usually not in danger. Fainting during exercise, without warning, or with palpitations or chest pain suggests a possible cardiac cause and must be evaluated by a doctor.}

\subsection{Blood Pressure Emergencies}

\subsubsection{High Blood Pressure Emergency (Hypertensive Emergency)}

Most people with high blood pressure (hypertension) feel no symptoms. Rarely, very high blood pressure causes acute damage to organs (brain, heart, eyes, kidneys) -- this is a hypertensive emergency and requires urgent hospital care.

\textbf{Recognize -- seek emergency care if these occur together or worsen:}
\begin{itemize}[leftmargin=*]
  \item A very severe headache, different from the person's usual headaches
  \item Blurred or double vision, or visual disturbances
  \item Chest pain or shortness of breath
  \item Confusion, personality changes, or difficulty speaking
  \item Nausea and vomiting, or weakness/numbness on one side of the body
  \item Nosebleeds that will not stop
\end{itemize}

\textbf{Management:}
\begin{enumerate}[leftmargin=*]
  \item Call emergency services immediately
  \item Have the person sit or lie down comfortably and remain calm
  \item Loosen tight clothing
  \item Do NOT give any blood pressure medication that is not theirs, and do not take extra doses of their own medication without medical direction
  \item Monitor responsiveness and breathing; be ready to begin CPR
  \item Do NOT ignore these symptoms as "just a headache" -- without treatment, a hypertensive emergency can cause stroke, heart attack, or organ damage
\end{enumerate}
\subsubsection{Low Blood Pressure (Hypotension)}

Low blood pressure (hypotension) is only an emergency when it causes symptoms, most commonly lightheadedness or fainting when standing up. It can follow dehydration, bleeding, infection, or allergic reactions.

\begin{itemize}[leftmargin=*]
  \item If the person feels faint: lay them flat, raise the legs only for non-traumatic causes, and let them recover gradually
  \item Give small sips of water if dehydration is suspected and the person is alert and able to swallow
  \item Seek emergency care if low blood pressure is accompanied by confusion, chest pain, shortness of breath, black/tarry stools, vomiting blood, or signs of severe allergic reaction
  \item A one-time reading on a home monitor with no symptoms is rarely an emergency -- but any unexplained, repeated, or severe readings should be discussed with a doctor
\end{itemize}

\subsection{Heat Illness}

See Chapter~\ref{ch:environmental} for detailed management of heat exhaustion and heat stroke.

\subsection{Poisoning}

Poisoning can occur through ingestion, inhalation, skin contact, or injection. Signs vary widely but may include nausea, vomiting, confusion, drowsiness, seizures, breathing difficulty, and burns around the mouth.

\textbf{General Principles:}

\begin{itemize}
  \item Remove the substance from the person if safe (brush off powder, flush chemicals with water) -- protect yourself with gloves or a barrier
  \item Do NOT induce vomiting unless instructed by poison control or medical professionals
  \item Call poison control center or emergency services immediately
  \item Keep the product container/poison source and any remains for medical personnel to identify
  \item Note the substance, approximate amount, and time of exposure
  \item If the person vomits, position them on their side to protect the airway
  \item Monitor breathing and level of consciousness; be ready to begin CPR
\end{itemize}

\textbf{Poison Control Numbers:}

\begin{center}
\small
\begin{tabularx}{\textwidth}{@{}>{\raggedright\arraybackslash}lX>{\raggedright\arraybackslash}r@{}}
\hline
\textbf{Location} & \textbf{Contact} & \textbf{Notes} \\
\hline
India (AIIMS NPIC) & 011-26589391, 011-26593677 & 24-hour helpline, New Delhi \\
India (State Poison Cells) & Check local health department & Regional centers in major cities \\
United States & 1-800-222-1222 & 24/7, American Association of Poison Control Centers \\
United Kingdom & 111 or 999 & NHS non-emergency or emergency \\
Australia & 13 11 26 & Poisons Information Centre \\
European Union & Local poison control center & Varies by country \\
\hline
\end{tabularx}
\end{center}

\textbf{Note}: Program the nearest poison control number into your phone before you need it. In India, the National Poisons Information Centre at AIIMS, New Delhi, is the primary reference for toxicological emergencies and can guide treatment decisions for common household poisons, snake antivenom selection, and pesticide exposures.

\subsubsection{Opioid Overdose}

Opioids (heroin, fentanyl, morphine, oxycodone, and similar) slow and can stop breathing. Overdose may occur with prescription painkillers or street drugs.

\begin{itemize}[leftmargin=*]
  \item \textbf{Recognize}: extreme drowsiness or unresponsiveness, slow or absent breathing, pinpoint pupils, blue-gray lips or nails
  \item Call emergency services immediately -- opioid overdose is reversible \cite{ilcor_2025,aha_arc_2024}
  \item Administer naloxone (Narcan) if available and you are trained: one dose into the nose or muscle, repeat every 2-3 minutes if no response
  \item Begin rescue breathing or CPR if breathing is absent but a pulse is present
  \item Naloxone wears off in 30-90 minutes -- continue to monitor and keep the person with medical personnel even if they wake up
  \item Do NOT leave the person alone; a second overdose can occur as naloxone wears off
\end{itemize}

\subsubsection{Alcohol Poisoning}

Alcohol can suppress the brain's control of breathing, heart rate, and gag reflex. Severe intoxication can be fatal.

\begin{itemize}[leftmargin=*]
  \item \textbf{Recognize}: confusion or stupor, vomiting, cold/clammy skin, slow or irregular breathing, low body temperature, seizures, unresponsiveness
  \item Call emergency services for anyone who is unconscious, has been vomiting repeatedly, or cannot be roused
  \item Turn an unconscious or vomiting person onto their side to prevent choking
  \item Do NOT give food, coffee, or more alcohol to "sober them up"
  \item Do NOT leave an intoxicated person to "sleep it off" -- their breathing can stop
\end{itemize}

\subsubsection{Carbon Monoxide (CO) Poisoning}

CO is an odorless, colorless gas from fires, vehicle exhaust, generators, and faulty heaters. It displaces oxygen from the blood.

\begin{itemize}[leftmargin=*]
  \item \textbf{Recognize}: headache, dizziness, nausea, weakness, confusion, chest pain, and cherry-red lips in severe cases (rare); multiple people becoming ill in the same building or vehicle is a clue
  \item Move the person to fresh air immediately -- do not enter an enclosed space without ventilation unless you can stay safe
  \item Call emergency services; anyone exposed needs medical evaluation even if they feel better in fresh air
  \item Do NOT assume recovery is complete -- CO poisoning requires oxygen therapy at a hospital
\end{itemize}

\noindent\textbf{\textcolor{red}{PREVENTION}}: Install carbon monoxide detectors near sleeping areas and test them monthly. Never run generators, cars, or heaters indoors or in attached garages.


\subsection{Caustic or Acid Ingestion}

Swallowing acids (cleaners, automotive fluids, battery acid) or other corrosive chemicals can cause severe internal burns to the mouth, throat, esophagus, and stomach. This is a medical emergency.

\textbf{Immediate Actions:}

\begin{enumerate}[leftmargin=*]
  \item Call poison control center (in India: National Poisons Information Centre, AIIMS New Delhi, 011-26589391; US: 1-800-222-1222) or emergency services immediately — provide the substance name and amount if known
  \item Do NOT induce vomiting — Vomiting can cause further damage as the corrosive material passes back up through the esophagus
  \item If the person is conscious, alert, and able to swallow safely, give small sips of water or milk to dilute the substance — only do this if advised by poison control or medical personnel
  \item If the substance is on the skin or in eyes, flush with copious running water for at least 15-20 minutes
  \item Keep the person calm and still — movement increases blood flow and potential absorption
  \item Save the product container or label for medical personnel to identify the substance
  \item NEVER give activated charcoal for caustic ingestions — it can obscure visibility during endoscopy and may cause harm
\end{enumerate}
\noindent\textbf{\textcolor{red}{WARNING}}: Signs of serious injury include drooling, difficulty swallowing, voice changes, abdominal pain, vomiting (especially with blood), or any alteration in consciousness. Seek immediate hospital care even if the person seems fine initially — internal damage can progress silently.

\subsection{Poisonous Plants (Poison Ivy, Poison Oak, Poison Sumac)}

The oil urushiol in these plants causes an itchy contact rash in most people. The rash is a reaction to the oil -- it does NOT spread from scratching, but the oil itself can spread if it remains on skin, clothing, or tools.

\begin{itemize}[leftmargin=*]
  \item \textbf{Prevention}: Learn to identify the plants ("leaves of three, let it be"); wear long sleeves, pants, and gloves in wooded areas; wash skin and clothing after suspected contact
  \item \textbf{Immediate action}: Wash exposed skin with soap and water as soon as possible (within 30 minutes); wash clothing and tools that touched the plant
  \item \textbf{Treat the rash}: wash the area gently; apply calamine lotion or hydrocortisone cream; cool compresses and oatmeal baths ease itching
  \item \textbf{Relief for itching}: oral antihistamines (diphenhydramine) may help, especially at night
  \item Do NOT scratch -- broken skin can become infected
\end{itemize}

\textbf{When to Seek Medical Care}: a severe or spreading rash, rash on the face, eyes, or genitals, blistering over large areas, difficulty breathing, or signs of infection (pus, spreading redness, fever). The rash may take 1--3 weeks to resolve.

\textbf{WARNING:} \textbf{Do not burn poison ivy, oak, or sumac.} Smoke from burning these plants carries urushiol and can cause severe, dangerous reactions in the lungs.

\subsection{Mushroom Poisoning}

Mushroom poisoning is difficult to identify and can be life-threatening. Many toxic mushrooms look like edible ones, so assume any wild mushroom that made someone ill is toxic.

\begin{enumerate}[leftmargin=*]
\item Call poison control (in India: National Poisons Information Centre, AIIMS New Delhi, 011-26589391) or emergency services immediately -- do not wait for symptoms to worsen
\item Keep samples: put a whole mushroom in a paper bag or take a clear photo for identification; do not touch it with bare hands
\item Do NOT induce vomiting unless specifically directed by poison control
\item Treat symptoms as they appear: for vomiting and diarrhea, replace fluids if the person can drink; for confusion, seizures, or difficulty breathing, call emergency services and monitor the airway
\item Some toxic mushrooms cause symptoms hours or days later -- even a mild reaction should be reported to poison control
\end{enumerate}

\textbf{WARNING:} \textbf{Never eat wild mushrooms unless positively identified by an expert.} "Deadly" mushrooms need not look dramatic -- the death cap mushroom is plain-looking and causes liver failure.

\subsection{Venomous Animal Bites (Snakes, Spiders, Scorpions)}

Venomous bites can cause local tissue damage, systemic effects, or both. Immediate action focuses on limiting venom spread and getting professional medical help.

\subsubsection{General Principles}

\begin{itemize}[leftmargin=*]
  \item Stay calm and keep the victim calm — fear and movement increase heart rate and venom spread
  \item Immobilize the bitten limb and keep it at or below heart level
  \item Remove constrictive items (rings, watches, tight clothing) near the bite site before swelling begins
  \item Note the time of the bite — this helps medical personnel track progression
  \item If safe, take a clear photo of the animal for identification — DO NOT attempt to capture or kill the animal
\end{itemize}

\subsubsection{What NOT to Do}

\begin{itemize}[leftmargin=*]
  \item DO NOT cut the wound or attempt to suck out venom — this causes additional tissue damage and infection risk without removing meaningful venom
  \item DO NOT apply ice to the bite site — ice can worsen tissue damage from certain venoms
  \item DO NOT use a tourniquet — modern first aid guidelines advise against tourniquets for snake bites; they can cause limb loss without reliably stopping venom spread
  \item DO NOT try to apply electric shocks or other folk remedies
\end{itemize}

\subsubsection{When to Escalate to Emergency Care}

Seek immediate emergency care if:
\begin{itemize}[leftmargin=*]
  \item The bite is from a known venomous species (in regions where these are common)
  \item Rapid swelling, pain, or discoloration develops around the bite
  \item Systemic symptoms appear: nausea/vomiting, dizziness, blurred vision, numbness, tingling, difficulty breathing, rapid heartbeat, or weakness
  \item The bite is on the face, neck, or torso (venom spreads faster to vital organs)
  \item The victim is a child, elderly, or has underlying health conditions
\end{itemize}

\noindent\textbf{Note}: In many regions, most "snake bites" are from non-venomous species. However, when in doubt, treat all bites as potentially venomous until a healthcare provider can assess.

\subsubsection{Snakebite in India}

India has a high burden of venomous snakebite. The \textbf{"Big Four"} medically significant species are the \textbf{Indian cobra} (Naja naja), \textbf{common krait} (Bungarus caeruleus), \textbf{Russell's viper} (Daboia russelii), and \textbf{saw-scaled viper} (Echis carinatus). Venomous bites occur predominantly in rural and agricultural areas, often at night or during the monsoon season, and are a leading cause of snakebite deaths worldwide.

\begin{itemize}[leftmargin=*]
  \item \textbf{Appearance}: Cobras have a hood and smooth scales; kraits are glossy with white cross-bands; Russell's viper and saw-scaled viper are stocky with keeled scales and triangular heads. Do not rely on appearance alone -- many harmless species mimic these features. A photo taken from a safe distance (do not approach or capture the snake) is the best aid to doctors for choosing antivenom
  \item \textbf{Typical symptoms}: Viper bites commonly cause painful swelling and bleeding (e.g., from the bite or gums); cobra and krait bites may show minimal local pain but can cause drooping eyelids, double vision, weakness, breathing difficulty, or paralysis within hours. A krait bite is often painless and may occur while the victim sleeps
  \item \textbf{First aid}: Follow the general principles above. Keep the victim calm and still, immobilize the limb, remove rings and tight clothing, and transport to the nearest hospital with antivenom facilities as quickly as possible. Do NOT apply a tourniquet, cut, or suck the wound
  \item \textbf{Do NOT use "snake stones", herbal preparations, or witch doctors' remedies} -- these cause dangerous delays in reaching definitive care
  \item \textbf{Antivenom}: Definitive treatment is hospital-based polyvalent antivenom. The treating doctor will decide based on signs of systemic envenoming. The national emergency number (112) and the National Poisons Information Centre (AIIMS New Delhi, 011-26589391) can help locate treatment facilities \cite{npic}
\end{itemize}

\subsubsection{Scorpion Stings}

In India, the \textbf{Indian red scorpion} (Hottentotta tamulus) causes severe envenoming, mainly in children, with intense local pain, sweating, and potentially life-threatening cardiovascular effects. First aid: keep the victim calm and still, apply a cold pack to reduce pain, and transport immediately to a hospital. Antivenom and supportive drugs are available in major centers. Note that in many states, 108/112 ambulance teams can be reached during the rainy season when scorpion stings peak.

\subsection{Animal and Human Bites (Non-Venomous)}

Dog, cat, and human bites are common and can become infected quickly. Cat bites are especially dangerous because their teeth create deep puncture wounds that seal over and trap bacteria.

\subsubsection{First Aid}

\begin{enumerate}[leftmargin=*]
\item Control bleeding with direct pressure and a clean dressing
\item Wash the wound thoroughly with mild soap and plenty of running water for several minutes
\item Apply an antibiotic ointment if available and cover with a sterile dressing
\item Change the dressing daily and watch for signs of infection (increasing redness, warmth, swelling, pus, red streaks, fever)
\end{enumerate}
\subsubsection{When to Seek Medical Care}

\begin{itemize}[leftmargin=*]
\item The bite broke the skin (especially cat or human bites)
\item Deep puncture wounds or heavy bleeding
\item Signs of infection develop
\item The bite is on the face, hands, or over a joint
\item Tetanus vaccination is unknown or more than 5 years old for a dirty wound
\item A dog, bat, raccoon, skunk, fox, or other potentially rabid animal caused the bite -- wash immediately and seek evaluation for rabies post-exposure treatment
\item The person is immunocompromised, diabetic, or elderly
\end{itemize}

\textbf{WARNING:} If the biting animal is wild, stray, or acting strangely, it may carry rabies. Do not attempt to capture it. Report the bite to local animal control and seek medical care promptly.

\subsection{Fever}

A fever is an elevated body temperature (typically over 100.4\textdegree F / 38\textdegree C) that helps fight infection. In adults, fever alone is rarely an emergency, but it can signal a serious underlying infection.

\subsubsection{Management}

\begin{itemize}[leftmargin=*]
\item Keep the person comfortable, cool, and well hydrated with small sips of fluids
\item Light clothing and a cool environment; a lukewarm sponge bath may help -- do NOT use cold water or ice (this causes shivering, which raises temperature)
\item Over-the-counter fever reducers (acetaminophen or ibuprofen) if appropriate and the person is not allergic
\item Do NOT give aspirin to children or teenagers (risk of Reye's syndrome)
\item Do NOT treat every fever -- a moderate fever is part of the immune response
\end{itemize}

\subsubsection{When to Seek Emergency Care}

\begin{itemize}[leftmargin=*]
\item Infant under 3 months with any temperature over 100.4\textdegree F (38\textdegree C) -- medical emergency
\item Fever with stiff neck, severe headache, confusion, seizures, or a rash that does not blanch when pressed
\item Difficulty breathing, persistent vomiting, or dehydration
\item Fever in an immunocompromised person or one with a serious chronic illness
\item Fever over 103\textdegree F (39.4\textdegree C) in an adult that does not respond to fever reducers
\item Any fever lasting more than 3 days
\end{itemize}

\section{Review Questions}

\begin{enumerate}[leftmargin=*]
\item List three symptoms of a heart attack that are more common in women.
\item Why must the exact time of stroke symptom onset be recorded?
\item A person has a seizure lasting 6 minutes. What should you do and why?
\item When should aspirin NOT be given to a person with suspected heart attack?
\item What are the key differences between managing a conscious and an unconscious person with hypoglycemia?
\item An unresponsive person has pinpoint pupils and very slow breathing. What condition should you suspect and what medication may help?
\end{enumerate}
\index{heart attack}
\index{myocardial infarction}
\index{stroke}
\index{FAST assessment}
\index{seizure management}
\index{epilepsy}
\index{hypoglycemia}
\index{hyperglycemia}
\index{diabetic emergency}
\index{anaphylaxis}
\index{allergic reaction}
\index{epinephrine}
\index{asthma attack}
\index{syncope}
\index{loss of consciousness}
\index{fainting}
\index{hypertensive emergency}
\index{high blood pressure}
\index{hypotension}
\index{low blood pressure}
\index{poisoning}
\index{opioid overdose}
\index{naloxone}
\index{alcohol poisoning}
\index{carbon monoxide}
\index{poison ivy}
\index{poison oak}
\index{poison sumac}
\index{urushiol}
\index{mushroom poisoning}
\index{toxic plants}
\index{animal bites}
\index{fever}
\index{chest pain}
\index{difficulty breathing}
